% Бібліографія (ДСТУ 8302:2015-подібне оформлення). Н1, Х3, G5 виключено (рішення Q7).
\begin{thebibliography}{99}
\footnotesize
\addcontentsline{toc}{chapter}{Бібліографія}
\bibitem{Singh2004} Singh R., Rogister A., Kaw P. The effect of asymmetric gas puffing
  on toroidal flow in the edge of tokamak plasma~// Physics of Plasmas. -- 2004. --
  Vol.~11, No.~1. -- P.~129--139. -- DOI: 10.1063/1.1633558. [C01]

\bibitem{Zou2013} Zou W., Li F., Lv B. On a nonlocal problem for a confined plasma in
  a Tokamak~// Applications of Mathematics. -- 2013. -- Vol.~58, No.~6. --
  P.~609--642. [C02]

\bibitem{Sigmar1987} Sigmar D.\,J., Zanino R., Hsu C.\,T. Evolution of poloidal
  variation of impurity density and ambipolar potential in rotating tokamak plasma.
  Part~I~: Report PFC/RR-87-8. -- MIT Plasma Fusion Center, 1987. -- 16~p. [C03]

\bibitem{Wong1987} Wong K.\,L., Cheng C.\,Z. Neoclassical diffusion of heavy
  impurities in a rotating tokamak plasma~: Report PPPL-2470. -- Princeton Plasma
  Physics Laboratory, 1987. -- 13~p. [C04]

\bibitem{Ward1992} Ward D.\,J., Jardin S.\,C. The effects of plasma deformability on
  the feedback stabilization of axisymmetric modes in tokamak plasmas~: Report
  PPPL-2805. -- Princeton Plasma Physics Laboratory, 1992. -- 59~p. [C05]

\bibitem{Hoppe2022} Hoppe M., Ekmark I., Berger E., Fülöp T. Runaway electron
  generation during tokamak start-up~// Journal of Plasma Physics. -- 2022. --
  Vol.~88. -- 905880317. -- DOI: 10.1017/S002237782200054X. [C06]

\bibitem{Mazzucato2000} Mazzucato E. The importance of the toroidal magnetic field
  for the feasibility of a tokamak burning plasma experiment~: preprint. -- Princeton
  Plasma Physics Laboratory, 2000. -- 9~p. [C07]

\bibitem{Fable2022} Fable E., Janky F., Treutterer W. et al. The modeling of a
  tokamak plasma discharge, from first principles to a flight simulator~// Plasma
  Physics and Controlled Fusion. -- 2022. -- Vol.~64. -- 044002. --
  DOI: 10.1088/1361-6587/ac466b. [C08]

\bibitem{Park1989} Park H.\,K., Goldston R.\,J., Taylor G. Study of electron density
  and pressure profile behavior in a tokamak plasma~: Report PPPL-2660. -- Princeton
  Plasma Physics Laboratory, 1989. -- 13~p. [C09]

\bibitem{Seto2014} Seto H., Fukuyama A. Formulation of two-dimensional transport
  modeling in tokamak plasmas~// Plasma and Fusion Research. -- 2014. -- Vol.~9. --
  1403002. -- DOI: 10.1585/pfr.9.1403002. [C10]

\bibitem{vanVugt2019} van Vugt D.\,C., Huijsmans G.\,T.\,A., Hoelzl M., Loarte A.
  Kinetic modeling of ELM-induced tungsten transport in a tokamak plasma~// Physics of
  Plasmas. -- 2019. -- Vol.~26, No.~4. -- 042508. -- DOI: 10.1063/1.5092319. [C11]

\bibitem{Ida2001} Ida K., CHS Group, LHD Group, JFT-2M Group. Structure of the radial
  electric field and toroidal/poloidal flow in high temperature toroidal plasma~//
  J.~Plasma Fusion Res. SERIES. -- 2001. -- Vol.~4. -- P.~36--42. [C12]

\bibitem{Hanada1990} Hanada K., Tanaka H., Iida M. et al. Sawtooth stabilization by
  localized electron cyclotron heating in a tokamak plasma~: Report PPPL-2725. --
  Princeton Plasma Physics Laboratory, 1990. -- 16~p. [C13]

\bibitem{AhoMantila2011} Aho-Mantila L. Divertor plasma conditions and their effect on
  carbon migration in the ASDEX Upgrade tokamak~: doctoral dissertation. -- Espoo~:
  Aalto University; VTT Publications 773, 2011. [C14]

\bibitem{Pigarov2006} Pigarov A.\,Yu., Smirnov R.\,D., Krasheninnikov S.\,I.,
  Rognlien T.\,D., Rozenberg M. Transport of dust particles in tokamak devices~:
  preprint UCRL-JRNL-221928. -- LLNL, 2006. -- 17~p. [C15]

\bibitem{Merlo2011} Merlo P. Plasma boundary reconstruction and shape control of
  tokamak discharges in RFX-mod~: tesi di laurea magistrale. -- Padova~: Università
  degli Studi di Padova, 2011. -- 170~p. [C16]

\bibitem{Hulse1992} Hulse R.\,A. Modeling of impurity transport in the core plasma~:
  Report PPPL-CFP-2776. -- Princeton Plasma Physics Laboratory, 1992. -- 28~p. [C17]

\bibitem{Jardin2000} Jardin S.\,C., Kaye S., Menard J., Kessel C., Glasser A.\,H.
  Tokamak Simulation Code modeling of NSTX~: Report PPPL-3472. -- Princeton Plasma
  Physics Laboratory, 2000. -- 7~p. [C18]

\bibitem{Fukuyama1992} Fukuyama A., Kasai T., Furutani Y. Transport simulasion [sic]
  in a burning tokamak plasma~// Memoirs of the Faculty of Engineering, Okayama
  University. -- 1992. -- Vol.~26, No.~2. -- P.~93--109. [C19]

\bibitem{Wang2006} Wang W.\,X., Lin Z., Tang W.\,M., Lee W.\,W. et al. Gyrokinetic
  simulation of global turbulent transport properties in tokamak experiments~: Report
  PPPL-4142. -- Princeton Plasma Physics Laboratory, 2006. -- 34~p. [C20]

\bibitem{Martsenyuk2025} Марценюк Ю.\,П. Характеристики газорозрядної плазми та
  розрядів різних типів у термоядерних дослідженнях~: дис. \ldots\ д-ра філософії~:
  104. -- Харків~: ННЦ ХФТІ НАН України, 2025. -- 150~с. [G1]

\bibitem{Marchuk2025} Марчук Ю.\,М. Моделювання електронів-утікачів, парних зіткнень
  та процесів у нижньому гібридному резонансі в плазмі з магнітним полем~: дис.
  \ldots\ д-ра філософії~: 104. -- Харків~: ІФП ННЦ ХФТІ НАН України, 2025. --
  147~с. [G2]

\bibitem{Romashchenko2021} Ромащенко О.\,В. Динаміка та фазові стани макрочастинок у
  пучково-плазмових системах~: автореф. дис. \ldots\ д-ра фіз.-мат. наук~: 01.04.08.
  -- Харків~: ХНУ ім. В.\,Н.~Каразіна, 2021. -- 36~с. [G3]

\bibitem{Tyshchenko2021} Тищенко М.\,Г. Поширення альфвенових хвиль та перенесення
  енергії поперек магнітних поверхонь у тороїдальній плазмі~: дис. \ldots\ канд.
  фіз.-мат. наук~: 01.04.08. -- Харків, 2021. -- 140~с. [G4]

\bibitem{Duba2018} Дуба І.\,Е. Сучасна цифрова багатозв'язна система управління
  основним контуром -- положення плазмового шнура -- установки термоядерного синтезу
  ТОКАМАК~// Вчені записки ТНУ ім. В.\,І.~Вернадського. Сер.: Технічні науки. --
  2018. -- Т.~29(68), №~4(1). -- С.~132--137. [Н2]

\bibitem{Moskvitina2010} Москвитина Ю.\,К. Численное моделирование влияния
  гофрировки на движение заряженной частицы в токамаке ITER~// Вісник ХНУ ім.
  В.\,Н.~Каразіна. Сер.: Математичне моделювання. -- 2010. -- №~925, вип.~14. --
  С.~147--154. [Н3]

\bibitem{Zaitsev2003} Зайцев Б.\,Ф., Асаёнок А.\,В. Напряженное состояние в винтовой
  обмотке тороидальной магнитной системы электрофизической установки под действием
  температуры~// Вестник НТУ «ХПИ». Динамика и прочность машин. -- 2003. -- №~12,
  т.~1. -- С.~64--71. [Х1]

\bibitem{Zaitsev2001} Зайцев Б.\,Ф., Шульженко Н.\,Г. Управление напряженным
  состоянием плоской многослойной катушки тороидального магнитного поля~// Вестник
  НТУ «ХПИ». Динамика и прочность машин. -- 2001. -- №~25. -- С.~86--90. [Х2]

\bibitem{Belyaeva2016} Беляева А.\,И., Галуза А.\,А., Коленов И.\,В., Савченко А.\,А.
  Роль рекристаллизации вольфрама в формировании шероховатости его поверхности под
  влиянием последовательного воздействия нейтронов и распыления~// Металлофизика и
  новейшие технологии. -- 2016. -- Т.~38, №~8. -- С.~1077--1102. [Х4]

\bibitem{ProjRegisters} Реєстри рішень проєкту «Магнітна рівновага, якої не видно»~:
  ESTABLISHED.md, CONJECTURES.md, open-questions.md~/ тека \texttt{Our try/00-decisions}. --
  Київ~: ФМФ КПІ, 2026. -- Редакція від 01.10.2026 (E1--E38, R1--R12; журнал змін CHANGELOG.md). [P1]

\bibitem{ProjDeepDives} Глибокі розбори D1--D3 та проби новизни~: код і результати~/
  теки \texttt{Our try/03-deep-dives}, \texttt{Our try/04-novelty}. -- 2026. [P2]

\bibitem{ProjGSsolver} Tokamak-GS-solver~: розв'язувач рівняння Ґреда--Шафранова з вільною межею
  (FreeGS-подібна ітерація Пікара) та PINN на даних KSTAR EFIT~: програмний код~/
  тека \texttt{Tokamak-GS-solver}. -- 2026. [P3]

\bibitem{ProjViz} tokamak-3d-viz~: візуалізація рівноваги DIII-D \#203702 (кадр 75, $t=1760$~мс),
  трасування силових ліній, острови, перерізи Пуанкаре~: програмний код (MIT)~/
  тека \texttt{tokamak-3d-viz}. -- 2026. Дані рівноваги: Fusion Equilibrium Challenge
  (Sophelio / General Atomics), CC BY 4.0. [P4]

\bibitem{FEC2026} Nakkina et al. The Fusion Equilibrium Challenge: Inferring Magnetic Geometry
  Without Magnetic Diagnostics~// arXiv:2609.01750. -- 2026. Набір даних
  \texttt{Sophelio/fusion-equilibrium-challenge}, випущено під CC BY 4.0; 1\,206 розрядів MAST походять із FAIR-MAST (CC BY-SA 4.0) --- правова підстава перевипуску не з'ясована (відкрите питання Q13 проєкту). [P5]

\bibitem{JET1975} The JET Project~: Design Proposal for the Joint European Torus~: EUR 5516e (EUR-JET-R5). --
  Luxembourg~: Commission of the European Communities, 1976. -- 678~p. -- © CECA, CEE, CEEA. [JET1975]

\bibitem{JETJU1983} JET Joint Undertaking. Progress Report 1983~: EUR 9472 EN (EUR-JET-PR1)~/ ed. B.\,E.~Keen,
  P.~Kupschus. -- Luxembourg~: ECSC/EEC/EURATOM, 1984. -- 148~p. [JET-PR1983]

\bibitem{JETBrochure1982} JET Joint Undertaking. Joint European Torus~: EUR 8306 EN (EUR-JET-RI~1). --
  Luxembourg~: ECSC/EEC/EURATOM, 1982. -- 46~p. [JET-RI1]

\bibitem{JETJU1979} JET Joint Undertaking. Annual Report 1979~: EUR-6832 EN (EUR-JET-AR2)~/ ed. B.\,J.~Green,
  G.\,W.~O'Hara. -- Luxembourg~: ECSC/EEC/EURATOM, 1980. -- 72~p. [JET-AR1979]

\bibitem{Wesson1999} Wesson~J. The Science of JET~: JET-R(99)13. -- Abingdon~: JET Joint Undertaking, 1999. -- 189~p. [Wesson]

\bibitem{ProjJET} tokamak-3d-viz~: 3D-реконструкція JET (проєкт 1975 і стан 1983) --- база розмірів
  \texttt{machines/jet1975}, \texttt{machines/jet1983}, геометричне ядро, рівновага Соловйова, гофрування,
  сцена Blender, веб-переглядач~: програмний код~/ тека \texttt{tokamak-3d-viz}. -- 2026. [P6]

\bibitem{ProjGuide} Путівник гіпотезами проєкту~: що ми припускали, що перевірили, що спростували~/
  \texttt{guide/main.pdf}. -- 2026. [P7]

\bibitem{Cerfon2010} Cerfon~A.\,J., Freidberg~J.\,P. «One size fits all» analytic solutions to the Grad--Shafranov
  equation~// Physics of Plasmas. -- 2010. -- Vol.~17, No.~3. -- 032502. -- DOI: 10.1063/1.3328818.
  (Поза корпусом; використано лише як метод у коді проєкту.)
\bibitem{Landreman2026} Landreman~M. Analytic toroidal 3D MHD equilibria and steady Euler flows with invariant
  surfaces~// arXiv:2609.26742 [physics.plasm-ph]. -- 2026. -- Версія~2 від 28.09.2026.
\end{thebibliography}
